% v8, May 27 2026, Claude Opus 4.7

\documentclass[11pt]{article}

\usepackage[margin=1in]{geometry}
\usepackage{amsmath, amssymb, amsthm}
\usepackage{mathtools}
\usepackage{xcolor}
\usepackage{hyperref}
\hypersetup{colorlinks=true,
            linkcolor=blue!50!black,
            citecolor=blue!50!black,
            urlcolor=blue!50!black}

%---- math macros ----
\newcommand{\R}{\mathbb{R}}
\newcommand{\sgn}{\operatorname{sgn}}
\newcommand{\rank}{\operatorname{rank}}
\newcommand{\mlrank}{\operatorname{mlrank}}
\newcommand{\bm}[1]{\mathbf{#1}}
\newcommand{\bQ}{\mathbf{Q}}
\newcommand{\bC}{\mathbf{C}}
\newcommand{\bA}{\mathbf{A}}
\newcommand{\bF}{\mathbf{F}}
\newcommand{\Dmat}{D}

%---- theorem environments ----
\newtheorem{theorem}{Theorem}
\newtheorem{lemma}{Lemma}
\newtheorem{proposition}{Proposition}
\newtheorem{corollary}{Corollary}

\title{Algebraic Relations on Determinantal Tensors}
\author{}
\date{May 27, 2026}

\begin{document}
\maketitle

\section*{Setup and Statement}

For $n\ge 5$, let $A^{(1)},\dots,A^{(n)}\in\R^{3\times4}$ be Zariski-generic.
For each ordered $4$-tuple $(\alpha,\beta,\gamma,\delta)\in[n]^4$ define
$Q^{(\alpha\beta\gamma\delta)}\in\R^{3\times3\times3\times3}$ by
\[
  Q^{(\alpha\beta\gamma\delta)}_{ijk\ell}
   \;=\; \det\!\bigl[A^{(\alpha)}(i,:);\,A^{(\beta)}(j,:);\,A^{(\gamma)}(k,:);\,A^{(\delta)}(\ell,:)\bigr],
   \qquad 1\le i,j,k,\ell\le 3,
\]
the $4\times4$ determinant of the matrix formed by stacking one row from each
of the four matrices.

We seek a polynomial map $\bF:\R^{81n^4}\to\R^{N}$ such that:
\begin{enumerate}\setlength\itemsep{2pt}
\item $\bF$ does not depend on $A^{(1)},\dots,A^{(n)}$;
\item the degrees of the coordinate functions of $\bF$ do not depend on $n$;
\item for $\lambda\in\R^{n\times n\times n\times n}$ with $\lambda_{\alpha\beta\gamma\delta}\ne0$
      precisely on the non-identical tuples,
      $\bF\bigl(\lambda_{\alpha\beta\gamma\delta}Q^{(\alpha\beta\gamma\delta)}\bigr)=0$
      \emph{iff} there exist $u,v,w,x\in(\R^{*})^n$ with
      $\lambda_{\alpha\beta\gamma\delta}=u_\alpha v_\beta w_\gamma x_\delta$ for
      all non-identical $(\alpha,\beta,\gamma,\delta)$.
\end{enumerate}

\medskip
\noindent\textbf{Theorem.} \emph{Such a map exists.} Assemble the tensors into
a single $3n\times3n\times3n\times3n$ block tensor $\bQ$ whose
$(\alpha,\beta,\gamma,\delta)$ block is $Q^{(\alpha\beta\gamma\delta)}$, and let
$\bF$ send the input to the $5\times5$ minors of the four
$3n\times27n^{3}$ mode flattenings of $\bQ$. Each coordinate of $\bF$ is a
$5\times5$ determinant, hence degree $5$ independent of $n$, and $\bF$
satisfies the rank-one characterization.

\medskip
\noindent The answer is \textbf{YES}.

\section*{The key identity: a Tucker decomposition of \texorpdfstring{$\bQ$}{Q}}

Index the modes of $\bQ$ by $[n]\times[3]$, so that the
$\bigl((\alpha,i),(\beta,j),(\gamma,k),(\delta,\ell)\bigr)$ entry of $\bQ$ is
$Q^{(\alpha\beta\gamma\delta)}_{ijk\ell}$.

\begin{lemma}\label{lem:tucker}
$\bQ$ admits the Tucker decomposition
\begin{equation}\label{eq:tucker}
  \bQ \;=\; \bC\times_1\bA\times_2\bA\times_3\bA\times_4\bA,
\end{equation}
where $\bA=[A^{(1)};\dots;A^{(n)}]\in\R^{3n\times4}$ is the vertical
concatenation of the data matrices, and $\bC\in\R^{4\times4\times4\times4}$ is
the sign tensor
\[
  C_{abcd}=\begin{cases}\sgn(abcd) & a,b,c,d\in[4]\text{ distinct},\\ 0 & \text{otherwise.}\end{cases}
\]
\end{lemma}

\begin{proof}
For $(\alpha,i),(\beta,j),(\gamma,k),(\delta,\ell)\in[n]\times[3]$, the
definition of the Tucker product gives
\[
  (\bC\times_1\bA\times_2\bA\times_3\bA\times_4\bA)_{(\alpha,i),(\beta,j),(\gamma,k),(\delta,\ell)}
  =\sum_{a,b,c,d\in[4]}C_{abcd}\,\bA_{(\alpha,i),a}\bA_{(\beta,j),b}\bA_{(\gamma,k),c}\bA_{(\delta,\ell),d}.
\]
Because $\bA$ is the vertical concatenation, $\bA_{(\alpha,i),a}=A^{(\alpha)}_{ia}$,
and the sum over distinct $a,b,c,d$ weighted by $\sgn(abcd)$ is exactly the
Leibniz expansion of the $4\times4$ determinant of the rows
$A^{(\alpha)}(i,:),A^{(\beta)}(j,:),A^{(\gamma)}(k,:),A^{(\delta)}(\ell,:)$.
This equals $Q^{(\alpha\beta\gamma\delta)}_{ijk\ell}=\bQ_{(\alpha,i),(\beta,j),(\gamma,k),(\delta,\ell)}$.
\end{proof}

\noindent Equation \eqref{eq:tucker} shows the multilinear rank of $\bQ$ is at
most $(4,4,4,4)$: each mode-$m$ flattening factors through the $4$-column matrix
$\bA$, so its column space (resp.\ row space) has dimension at most $4$. In
particular \textbf{every $5\times5$ minor of every flattening of $\bQ$ vanishes},
and the same holds for any tensor of multilinear rank $\le(4,4,4,4)$. This is
why $\bF$, the collection of $5\times5$ minors of the four flattenings, is the
right object: it tests membership in the variety of multilinear rank
$\le(4,4,4,4)$ tensors.

\section*{Definition of \texorpdfstring{$\bF$}{F}}

Given input tensors $\{T^{(\alpha\beta\gamma\delta)}\in\R^{3\times3\times3\times3}\}$,
assemble the $3n\times3n\times3n\times3n$ block tensor $\mathbf T$ with blocks
$T^{(\alpha\beta\gamma\delta)}$. Let $\mathbf T_{(m)}\in\R^{3n\times27n^{3}}$ be
its mode-$m$ flattening, $m=1,2,3,4$. Define
\[
  \bF(\{T^{(\alpha\beta\gamma\delta)}\})
   \;:=\;\bigl(\text{all }5\times5\text{ minors of }\mathbf T_{(1)},\mathbf T_{(2)},\mathbf T_{(3)},\mathbf T_{(4)}\bigr).
\]
Each coordinate is a $5\times5$ determinant in the entries of the input, hence a
degree-$5$ polynomial in those entries, independent of $n$; and $\bF$ visibly
does not reference $A^{(1)},\dots,A^{(n)}$. Properties (1) and (2) hold. We
write $\lambda\odot_b\bQ$ for the blockwise scaling whose
$(\alpha,\beta,\gamma,\delta)$ block is $\lambda_{\alpha\beta\gamma\delta}Q^{(\alpha\beta\gamma\delta)}$.
Property (3) is the content of the next two sections.

\section*{The ``if'' direction}

\begin{proposition}\label{prop:if}
If $\lambda_{\alpha\beta\gamma\delta}=u_\alpha v_\beta w_\gamma x_\delta$ on the
non-identical tuples (and $0$ on the diagonal $\alpha=\beta=\gamma=\delta$), then
$\bF(\lambda\odot_b\bQ)=0$.
\end{proposition}

\begin{proof}
First, the diagonal blocks of $\bQ$ vanish: $Q^{(\alpha\alpha\alpha\alpha)}_{ijk\ell}$
is a determinant with rows drawn from the single $3\times4$ matrix $A^{(\alpha)}$,
and any $4\times4$ determinant of vectors lying in a $3$-dimensional row space is
$0$. Hence the value of $\lambda$ on the diagonal is irrelevant and
$\lambda\odot_b\bQ=(u\otimes v\otimes w\otimes x)\odot_b\bQ$.

Blockwise scaling by a rank-one tensor with nonzero entries is the same as a
Tucker multiplication by invertible diagonal matrices:
\[
  (u\otimes v\otimes w\otimes x)\odot_b\bQ
   \;=\;\bQ\times_1 D_u\times_2 D_v\times_3 D_w\times_4 D_x,
\]
where $D_u\in\R^{3n\times3n}$ is the diagonal matrix repeating each $u_\alpha$
three times (once per row of the $\alpha$ block), and similarly $D_v,D_w,D_x$.
Multiplying \eqref{eq:tucker} on each mode by an invertible matrix preserves the
multilinear rank, so $\lambda\odot_b\bQ$ still has multilinear rank
$\le(4,4,4,4)$, and all its $5\times5$ flattening minors vanish. Thus
$\bF(\lambda\odot_b\bQ)=0$.
\end{proof}

\section*{The ``only if'' direction}

Assume $\lambda$ has nonzero entries exactly off the diagonal and
$\bF(\lambda\odot_b\bQ)=0$, i.e.\ $\lambda\odot_b\bQ$ has multilinear rank
$\le(4,4,4,4)$. We show $\lambda$ is rank-one off the diagonal.

\subsection*{Normalization}

Replacing $\lambda$ by its entrywise product with $\bar u\otimes\bar v\otimes\bar w\otimes\bar x$,
where e.g.\ $\bar u_1=1$ and $\bar u_\alpha=\lambda_{\alpha111}^{-1}$ for
$\alpha\ge2$ (and $\bar v,\bar w,\bar x$ defined symmetrically on the other
modes), preserves both the multilinear rank of $\lambda\odot_b\bQ$ and the
property of being rank-one off-diagonal. Hence we may assume
\begin{equation}\label{eq:normalize}
  \lambda_{\alpha111}=\lambda_{1\beta11}=\lambda_{11\gamma1}=\lambda_{111\delta}=1
  \qquad\text{for all }\alpha,\beta,\gamma,\delta\in\{2,\dots,n\}.
\end{equation}
We will show there is a constant $c\in\R^{*}$ with
\begin{equation}\label{eq:goal}
  \lambda_{\alpha\beta\gamma\delta}=c^{\,m-1},
  \qquad m:=\#\{\text{indices among }\alpha,\beta,\gamma,\delta\text{ that are }\ne 1\},
\end{equation}
for all non-identical tuples. Then $u=v=w=(1,c,\dots,c)$ and
$x=(c^{-1},1,\dots,1)$ give $\lambda_{\alpha\beta\gamma\delta}=u_\alpha v_\beta w_\gamma x_\delta$,
establishing rank one.

\subsection*{Reading off $5\times5$ submatrix determinants}

The multilinear-rank hypothesis says every $5\times5$ minor of every flattening
of $\lambda\odot_b\bQ$ vanishes. We extract three families of such minors. Write
$(\lambda\odot_b\bQ)_{(1)}\in\R^{3n\times27n^3}$ for the mode-$1$ flattening;
rows are indexed by $(\alpha,i)\in[n]\times[3]$ and columns by
$((\beta,j),(\gamma,k),(\delta,\ell))$.

\medskip
\noindent\textbf{Step 1 (two non-unit indices).}
Choose the $5\times5$ submatrix with rows $(1,1),(1,2),(1,3),(\alpha,1),(\alpha,2)$
and a fixed list of five columns chosen so that, viewed as a polynomial in the
two variables $\lambda_{\alpha\alpha11}$ and $\lambda_{\alpha1\beta1}$, the
determinant has degree at most $2$; it vanishes when $\lambda_{\alpha1\beta1}=0$,
and also when $\lambda_{\alpha1\beta1}=\lambda_{\alpha\alpha11}$ (in the latter
case the submatrix becomes a submatrix of the unscaled $\bQ_{(1)}$ after row
operations, hence rank-deficient by Lemma~\ref{lem:tucker}). Therefore this
determinant equals
\[
  s\cdot\lambda_{\alpha1\beta1}\bigl(\lambda_{\alpha1\beta1}-\lambda_{\alpha\alpha11}\bigr),
\]
with $s=s(A^{(1)},A^{(\alpha)},A^{(\beta)})$ a polynomial in the data. A single
generic numerical instance (it suffices to take $\alpha=\beta$, the most
specialized case) shows $s\not\equiv0$, so $s\ne0$ Zariski-generically. The
minor vanishes, forcing $\lambda_{\alpha1\beta1}=\lambda_{\alpha\alpha11}$.
Applying the same argument to all modewise permutations $\pi$ (using
$\pi\cdot\bQ=\sgn(\pi)\bQ$, so the permuted flattening has the same rank bound)
gives $\lambda_{\pi(\alpha1\beta1)}=\lambda_{\pi(\alpha\alpha11)}$. Combining the
specializations $\alpha=\beta$ and $\alpha\ne\beta$ yields a single constant
\[
  c:=\lambda_{\pi(\alpha\beta11)}\quad\text{for all }\alpha,\beta\in\{2,\dots,n\}\text{ and all }\pi,
\]
i.e.\ every $\lambda$-entry with exactly two non-unit indices equals $c$.

\medskip
\noindent\textbf{Step 2 (one non-unit index).}
Take the $5\times5$ submatrix with rows $(1,1),(1,2),(1,3),(\alpha,1),(\alpha,2)$
and columns chosen so the determinant, viewed as a polynomial in $c$ and
$\lambda_{\alpha\beta\gamma1}$, has degree $\le3$, vanishes when $c=0$ (two rows
become dependent), and vanishes when $c^{2}=\lambda_{\alpha\beta\gamma1}$ (the
submatrix reduces to one of $\bQ_{(1)}$ after row/column operations). Hence the
determinant is a data-dependent scalar times $c\,(c^{2}-\lambda_{\alpha\beta\gamma1})$.
Generic non-vanishing of the scalar (checked on one instance with
$\alpha=\beta=\gamma$) and the rank hypothesis give
$\lambda_{\alpha\beta\gamma1}=c^{2}$, and by symmetry
$\lambda_{\pi(\alpha\beta\gamma1)}=c^{2}$ for all $\pi$. Thus every $\lambda$-entry
with exactly one non-unit index equals $c^{2}$.

\medskip
\noindent\textbf{Step 3 (no non-unit index).}
Take the $5\times5$ submatrix with rows $(1,1),(1,2),(1,3),(\alpha,1),(\alpha,2)$
and columns producing a determinant equal to a data scalar times
$c\,(c^{3}-\lambda_{\alpha\beta\gamma\delta})$, where now
$\alpha,\beta,\gamma,\delta\in\{2,\dots,n\}$ are not all equal (the construction
parallels Steps 1–2; the most specialized instance $\alpha=\beta=\gamma$ certifies
non-vanishing). The rank hypothesis forces $\lambda_{\alpha\beta\gamma\delta}=c^{3}$,
and by symmetry $\lambda_{\pi(\alpha\beta\gamma\delta)}=c^{3}$.

\medskip
Steps 1–3 establish \eqref{eq:goal}: a $\lambda$-entry with $m$ non-unit indices
equals $c^{\,m-1}$ (using the normalization \eqref{eq:normalize} for $m=1$). As
noted, this is exactly the rank-one factorization
$\lambda_{\alpha\beta\gamma\delta}=u_\alpha v_\beta w_\gamma x_\delta$ with
$u=v=w=(1,c,\dots,c)$ and $x=(c^{-1},1,\dots,1)$. This proves the ``only if''
direction and completes the proof of the Theorem.\qed

\section*{Degree, dimension, and independence}

\textbf{Degree.} Each coordinate of $\bF$ is a $5\times5$ determinant in the
entries of the input, so has degree $5$, independent of $n$. \checkmark

\textbf{Independence of the data.} $\bF$ is defined purely as minors of the
flattenings of the input block tensor; it never refers to $A^{(1)},\dots,A^{(n)}$. \checkmark

\textbf{Number of coordinates.} $N$ is the number of $5\times5$ minors of four
$3n\times27n^{3}$ matrices, which grows polynomially in $n$; only the
\emph{degree} is required to be bounded, which it is. \checkmark

\section*{Remarks}

\begin{enumerate}\setlength\itemsep{4pt}
\item \textbf{Why $5\times5$ and not $2\times2$.} The obstruction is multilinear
rank $(4,4,4,4)$: a flattening has rank $\le4$, so the vanishing of all
$5\times5$ minors is the exact algebraic condition. Smaller minors do not
characterize the variety, and larger ones are redundant.
\item \textbf{Role of genericity.} Genericity of the $A^{(i)}$ enters only in the
``only if'' direction, to certify that the data-dependent scalars $s$ in
Steps 1–3 are nonzero; each is verified by exhibiting one numerical instance,
after which Zariski-genericity does the rest.
\item \textbf{Diagonal blocks vanish.} $Q^{(\alpha\alpha\alpha\alpha)}=0$ because
four rows of a $3\times4$ matrix are linearly dependent; this is what lets the
off-diagonal rank-one structure of $\lambda$ be detected without the diagonal
interfering.
\item \textbf{Fourth-order analogue.} The construction is the order-$4$ analogue
of the corresponding third-order determinantal-variety result; the block-tensor
flattening and the multilinear-rank bound are the common mechanism.
\end{enumerate}

\begin{thebibliography}{9}
\bibitem{Landsberg} J.~M.~Landsberg.
  \emph{Tensors: Geometry and Applications}. AMS, 2012.
\bibitem{KoldaBader09} T.~G.~Kolda and B.~W.~Bader.
  \emph{Tensor decompositions and applications}.
  SIAM Review \textbf{51} (2009), 455--500.
\bibitem{Sturmfels} B.~Sturmfels.
  \emph{Algorithms in Invariant Theory}, 2nd ed., Springer, 2008.
\end{thebibliography}

\end{document}